\section{Fisher-Rao Quantization-Aware Distance and Local TurboQuant (C1)}
\label{sec:quantization}

This section presents the paper's hero contribution: a new distance metric for mixed-precision embeddings (FRQAD) and the first application of near-optimal data-oblivious quantization to persistent agent memory stores (LT2E).

\subsection{From KV Cache to Persistent Memory}

TurboQuant~\cite{turboquant} was designed for KV cache---vectors generated during inference, used briefly, and discarded. Persistent agent memory stores differ fundamentally: vectors have lifetimes of months, require random-access similarity search, grow unboundedly, and cannot be regenerated if corrupted.

We adopt TurboQuant's per-coordinate scalar quantization and extend it with cognitive lifecycle management.

\subsubsection{TurboQuant Algorithm}

\textbf{One-time setup:}
\begin{enumerate}[topsep=2pt, itemsep=1pt]
  \item Generate a random orthogonal rotation matrix $\mathbf{\Pi} \in \mathbb{R}^{d \times d}$ via QR decomposition of a Haar-distributed random Gaussian matrix~\cite{mezzadri}. Store on disk; reuse for all embeddings.
  \item Pre-compute Lloyd-Max optimal codebook centroids $c_1, \ldots, c_{2^b}$ for the Beta$(1/2, (d{-}1)/2)$ distribution on $[-1, 1]$, for each supported bit-width $b \in \{2, 4, 8\}$.
\end{enumerate}

\textbf{Quantize($\mathbf{x}$, $b$):} Rotate $\mathbf{y} = \mathbf{\Pi} \cdot \mathbf{x}$; for each coordinate $j$: assign $\text{idx}_j = \arg\min_k |y_j - c_k|$; pack indices into $b \cdot d$ bits.

\textbf{Dequantize:} Reconstruct $\tilde{y}_j = c_{\text{idx}_j}$; rotate back $\tilde{\mathbf{x}} = \mathbf{\Pi}^T \cdot \tilde{\mathbf{y}}$.

The key insight: after random orthogonal rotation, each coordinate follows a Beta$(1/2, (d{-}1)/2)$ distribution, converging to $\mathcal{N}(0, 1/d)$ in high dimensions. Per-coordinate scalar quantization is then near-optimal.

\begin{theorem}[MSE Distortion Upper Bound~{\cite{turboquant}}]
\label{thm:turboquant-mse}
For $b$-bit TurboQuant applied to any unit-norm vector $\mathbf{x} \in \mathbb{R}^d$:
\begin{equation}
D_{\text{mse}} \leq \sqrt{\frac{3\pi}{2}} \cdot \frac{1}{4^b}
\end{equation}
This is within $2.7\times$ of the information-theoretic lower bound.
\end{theorem}

\subsubsection{LT2E: Adaptations for Persistent Stores}

Three key modifications: (1)~\emph{Pre-computed rotation matrices} stored on disk and reused for all embeddings, eliminating per-embedding rotation cost. (2)~\emph{Mixed-precision storage}: 2/4/8/32-bit precision per embedding, selected by the EAP scheduler (Section~\ref{sec:forgetting}). (3)~\emph{Backward-compatible search} across precision levels via RRF with precision-aware weighting.

\subsection{Fisher-Rao Quantization-Aware Distance (FRQAD)}
\label{sec:frqad}

FRQAD is a new distance metric for comparing embeddings at different quantization levels. Our systematic literature search found \textbf{zero prior work} combining information geometry with vector quantization for similarity retrieval.

\textbf{Intuition.} Quantization introduces known noise. A 4-bit embedding has more uncertainty than a 32-bit original. Cosine similarity ignores this. FRQAD accounts for it by treating embeddings as parameters of probability distributions where quantization error determines variance.

\begin{definition}[FRQAD]
For memories $m_i, m_j$ with embeddings $\theta_i, \theta_j$ at bit-widths $b_i, b_j$, define:
\begin{equation}
d_{\text{FRQAD}}(m_i, m_j) = d_{\text{FR}}\left(\mathcal{N}(\theta_i, \sigma^2_{\text{eff},i} \mathbf{I}), \, \mathcal{N}(\theta_j, \sigma^2_{\text{eff},j} \mathbf{I})\right)
\end{equation}
where $\sigma^2_{\text{eff},k} = \sigma^2_{\text{obs}}(m_k) \cdot (32/b_k)^{\kappa}$ inflates the base observation variance by the quantization factor, and $d_{\text{FR}}$ is the Fisher-Rao geodesic on the Gaussian manifold~\cite{amari-info-geo}.
\end{definition}

For diagonal Gaussians, the full Atkinson-Mitchell geodesic~\cite{atkinson-mitchell}:
\begin{equation}
d_{\text{FR}} = \sqrt{\sum_{k=1}^{d} \left[\sqrt{2} \cdot \text{arccosh}\left(1 + \frac{(\mu_{1k} - \mu_{2k})^2 + 2(\sigma_{1k} - \sigma_{2k})^2}{4\sigma_{1k}\sigma_{2k}}\right)\right]^2}
\end{equation}

\textbf{Design note: simplified vs.\ full geodesic.} The SLM retrieval pipeline uses a \emph{graduated ramp} from cosine similarity to the simplified Fisher-Rao form ($d^2 = \sum(\mu_1 - \mu_2)^2/\sigma^2$) as facts accumulate access history. This simplified form is correct for uniform-precision retrieval (all float32), where all variances are equal. However, in mixed-precision scenarios---where some embeddings are quantized to 4-bit or 2-bit---inflating $\sigma^2$ in the denominator \emph{reduces} distance, producing incorrect rankings. FRQAD addresses this by using the full Atkinson-Mitchell geodesic with the variance-mismatch term. In the current release, all embeddings are stored at float32 by default, making FRQAD and standard Fisher-Rao produce identical results; FRQAD's advantage emerges when the EAP scheduler promotes mixed-precision storage.

\begin{proposition}[Monotonic degradation]
For fixed embeddings and observation variances, FRQAD distance increases monotonically as either embedding's bit-width decreases:
$b_i' < b_i \implies d_{\text{FRQAD}}(m_i|_{b_i'}, m_j) > d_{\text{FRQAD}}(m_i|_{b_i}, m_j)$
\end{proposition}

This guarantees that quantized (faded) memories are ranked lower than full-precision (active) memories without any explicit re-weighting.

% ============================================================================
% 5. THE FORGETTING BRAIN (C2)
% ============================================================================
